% Source: https://artofproblemsolving.com/community/c3h1772259_math_problems
% Stable identity: aops-topic-1772259-q04
Let $\angle FDB=y$. Since $DF=BF$, triangle $BDF$ is isosceles and $\angle FBD=y$, so its exterior angle at $F$ is $\angle BFE=2y$. Since $EF=BE$, triangle $BEF$ is isosceles with $\angle EBF=\angle EFB=2y$. Now $\angle DBC=45^\circ$ (diagonal of a square) and $\angle DBC=\angle DBF+\angle FBE=y+2y$, so $y=15^\circ$.

$F$ is equidistant from $B$ and $D$, so it lies on the perpendicular bisector of $BD$, which is the other diagonal $AC$. Therefore $\angle DAF=\angle DAC=45^\circ$. Also $\angle ADF=\angle ADB+\angle BDF=45^\circ+15^\circ=60^\circ$. In triangle $ADF$,
\[
\angle DFA=180^\circ-45^\circ-60^\circ=\boxed{75^\circ}.
\]

\begin{asy}
usepackage("amsmath");
size(7cm);
defaultpen(fontsize(10pt));
pair A=(0,1),B=(1,1),C=(1,0),D=(0,0);
pair Ep=(1,1/sqrt(3)); real x=sqrt(3)/(sqrt(3)+1); pair F=(x,x/sqrt(3));
draw(A--B--C--D--cycle,linewidth(0.9));
draw(D--Ep,linewidth(0.8)); draw(B--F,linewidth(0.8)); draw(A--F,red+linewidth(0.8));
draw(B--D,dashed+gray); draw(F--C,dashed+gray);
void tk(pair P,pair Q,int n){pair m=(P+Q)/2,u=unit(Q-P),v=(-u.y,u.x); for(int i=0;i<n;++i){pair c=m+(i-(n-1)/2)*0.03*u; draw(c-0.035*v--c+0.035*v,linewidth(0.7));}}
tk(D,F,1); tk(B,F,1); tk(F,Ep,2); tk(B,Ep,2);
draw(arc(F,0.07,degrees(A-F),degrees(D-F)),red); label("$75^\circ$",F+0.135*dir(182),red);
draw(arc(D,0.3,degrees(Ep-D),45),blue); label("$15^\circ$",D+0.4*dir(37.5),blue);
label("$A$",A,NW);label("$B$",B,NE);label("$C$",C,SE);label("$D$",D,SW);label("$E$",Ep,E);label("$F$",F,dir(-80));
label("\mbox{}",(min(currentpicture,user=true)+max(currentpicture,user=true))/2,black); // keeps the text after the figure black when LaTeX typesets the labels inline
\end{asy}

(The figure is to scale: $\angle EDC=30^\circ$, so $CE=\dfrac{1}{\sqrt3}$ for a unit square.)
